A group $G$ is said to be commutative if $x y = y x$.
A ring $K$ is said to be commutative (or abelian) if $x y = y x$.
$x$ is said to be a divisor of $y$ if $y = x z$.
$G$ and $H$ are said to be isomorphic if $f(G) = H$.
$x$ is called a unit if $x y = 1$.
An isomorphism is by definition any homomorphism that is bijective.
$p$ is called the characteristic of $K$.
The order of $G$ is defined to be the determinant of $G$.
We define the order of $G$ to be $n$.
We say that $G$ is abelian if $G$ is commutative.
We denote by $N$ the kernel of $f$.
The characteristic of $K$ is denoted by $c$.
Groups that are commutative are called abelian groups.
$f$ is a bijective homomorphism.
$G$ is not a commutative group.
$G$ has at least two subgroups.
$x$ has a left reflection.
For $x$ to be invertible it is necessary and sufficient that $x$ is a unit.
For an element $x$ of $G$ to be reflexible, it is necessary and sufficient that $x$ has a left reflection.
For $G$ to be commutative it is sufficient that $G$ is finite.
There exists a unique element $x$ of $G$ such that $x$ is invertible.
If $x$ is reflexible, then so is $y$.
$x$ is finite whenever $y$ is finite.
$f$ is injective (resp. surjective).
Every subgroup of $G$ that contains $x$ is commutative.
Every divisor of $x$ divisible by $y$ is finite.
Given an element $x$ of $G$, $x$ is reflexible.
The number of subgroups of $G$ is finite.
The set of all subgroups of $G$ is finite.
The kernel of $f$, called the order of $f$, is finite.
$x$ is equal to the unique element $y$ of $G$ such that $y$ is invertible.
A subgroup of $G$, namely $H$, is finite.
Every element $x \in G$ such that $x$ is invertible is reflexible.
Every element of $G$ of the form $x y$, where $y \in H$, is reflexible.
The following conditions are equivalent: a) $G$ is commutative; b) $G$ is abelian; c) $G$ is finite.
$G$ has exactly $n$ subgroups.
Given an element $x \in G$, $x$ is reflexible.
Let $H$ be a subgroup of $G$ with the following properties: a) $H$ is finite; b) $H$ is commutative. Then $H$ is normal.
